\documentclass[../../main.tex]{subfiles}
\begin{document}

{\bf [解答]}

与えられた条件(2)の
\begin{align}
  b\int_a^x f(t)dt=xf(x) \label{eq:1}
\end{align}
について，$x=a$とすると
\begin{align}
  a f(a) = 0 \label{eq:2}
\end{align}
である．

また条件(1)から\cref{eq:1}の両辺は $x$ で微分できて，
\begin{align}
  & b f(x) = f(x) + x f'(x)
  \therefore
  & (b-1) f(x) = x f'(x) \label{eq:3}
\end{align}
である．$y = f(x)$ として代入すると，
\begin{align}
  (b-1) y = x \dfrac{dy}{dx} \label{eq:4}
\end{align}
である．$y \neq 0$ とすると
\begin{align}
  \dfrac{b-1}{x} dx = \dfrac{1}{y} dy
\end{align}
だから両辺積分して，$C_0$ を定数として，
\begin{align}
  & (b-1) \log x = \log y + C_0
  \therefore
  & y = C \cdot x^{b-1} \quad (C\text{：定数})
\end{align}
である．$y=0$の時も$C=0$とすることでこの表式で良い．
\cref{eq:2}から
\begin{align}
  a \cdot C \cdot a^{b-1} = 0 \iff C = 0 \text{ or } a = 0
\end{align}
に注意すると，$C$ を定数として，
\begin{align}
  y = C x^{b-1}
\end{align}
である．

{\bf [解注]}

(1)に注意してさらにこまかく分類すれば，
\begin{align}
  \begin{cases}
    a=0 \wedge 1 \le b \text{ の時 } & y = C x^{b-1} \\
    a \neq 0 \text{ or } b < 1 \text{ の時 } & y = 0
  \end{cases}
\end{align}
となる．

\end{document}
